%% kapa: exam=1 round=36 year=2025 period=2 form=A subject=accounting subject_ko=회계학 seq=11 number=51
%% source: Q-Net 감정평가사 기출문제 (cst00309) artlSeq=5231525 file=raw/1008/5231525/2025년 제36회 감정평가사 1차시험 2교시 문제지 원본.pdf
%% conversion: kordoc 4.24.51 -> tools/gen1.py (auto); review=proofread-sonnet; flags=none
\begin{question}{51}
(주)감평은 20×1년 초 영업용 차량운반구(취득원가 ￦500,000, 내용연수 5년, 잔존가치 ￦0, 정액법 상각)를 취득하고 원가모형을 적용하였다. (주)감평은 20×2년 초 차량운반구의 일상적인 유지와 관련하여 ￦30,000을 지출하였다. 또한 동 일자에 차량운반구의 성능을 향상시키기 위하여 ￦200,000을 추가 지출하였고, 이로 인해 차량운반구의 내용연수는 2년 연장되었으며 잔존가치는 ￦60,000으로 증가되었다. 동 차량운반구와 관련된 회계처리로 인한 (주)감평의 20×2년 당기순이익 감소액은?
\begin{choices}
  \item ￦90,000
  \item ￦95,000
  \item ￦100,000
  \item ￦120,000
  \item ￦130,000
\end{choices}
\end{question}
\answer{4}
